\Needspace{11\baselineskip}
\begin{center}
\textit{EXKURS MIT EINER BEHANDLUNG DER TERNÄREN FORMEN.}
\end{center}

\Needspace{8\baselineskip}
\Article{266}
Bisher haben wir unsere Untersuchung auf solche Funktionen zweiten Grades beschränkt, die \textit{zwei} Unbestimmte enthalten, und mussten ihnen keine besondere Bezeichnung geben. Doch lässt sich dieser Gegenstand offenbar als ein sehr spezieller Teil der allgemeinsten Untersuchung \textit{über ganzrationale algebraische h\,o\,m\,o\,g\,e\,n\,e Funktionen mehrerer Unbestimmter und verschiedener Grade} auffassen. Diese Funktionen können wir zweckmäßig nach ihrem Grad als \textit{Formen zweiten, dritten, vierten Grades usw.}, nach der Anzahl ihrer Unbestimmten dagegen als \textit{binäre, ternäre, quaternäre Formen} usw.\ unterscheiden. Die Formen, die wir bisher einfach so genannt haben, werden also \textit{binäre Formen zweiten Grades} heißen; Funktionen wie
\[
Axx+2Bxy+Cyy+2Dxz+2Eyz+Fzz
\]
wobei $A$, $B$, $C$, $D$, $E$, $F$ gegebene ganze Zahlen bezeichnen) werden dagegen \textit{ternäre Formen zweiten Grades} heißen, und so weiter. Dieser Abschnitt ist zunächst allein den binären Formen zweiten Grades gewidmet. Doch sind noch mehrere Wahrheiten über sie übrig, und zwar besonders schöne, deren eigentliche Quelle in der Theorie der ternären Formen zweiten Grades zu suchen ist. Deshalb schieben wir hier einen kurzen Exkurs zu dieser Theorie ein und werden aus ihren ersten Grundlagen das behandeln, was zur Vervollständigung der Theorie der binären Formen nötig ist. Wir \SourceLeaf{305}hoffen, dass dies den Mathematikern willkommener sein wird, als wenn wir jene Wahrheiten entweder unterdrückten oder durch weniger natürliche Methoden gewännen. Eine eingehendere Untersuchung dieses sehr wichtigen Gegenstands müssen wir uns für eine andere Gelegenheit vorbehalten: zum einen, weil seine Fülle schon jetzt die Grenzen dieses Werks weit überschreiten würde, zum anderen, weil wir hoffen dürfen, dass er künftig noch durch bedeutende Fortschritte bereichert wird. Sowohl die quaternären, quinären usw.\ Formen zweiten Grades als auch alle Formen höherer Grade schließen wir hier gänzlich aus unserem Vorhaben aus\HistoricalNote{*)}{Deshalb sind im Folgenden immer binäre oder ternäre Formen \textit{zweiten Grades} gemeint, wenn wir einfach von solchen Formen sprechen.}. Es mag genügen, die Mathematiker auf dieses weite Feld aufmerksam zu machen, in dem sie überreiches Material finden werden, um ihre Kräfte zu erproben und die höhere Arithmetik durch hervorragende Fortschritte zu bereichern.

\Needspace{8\baselineskip}
\Article{267}
Es wird sehr zur Klarheit beitragen, unter den drei Unbestimmten einer ternären Form ebenso wie bei den binären Formen eine feste Reihenfolge einzuführen, sodass die \textit{erste, zweite} und \textit{dritte Unbestimmte} voneinander unterschieden werden. Bei der Anordnung der einzelnen Glieder der Form werden wir stets folgende Reihenfolge einhalten: An erster Stelle steht das Glied mit dem Quadrat der ersten Unbestimmten; danach folgen nacheinander die Glieder mit dem Quadrat der zweiten, dem Quadrat der dritten, dem doppelten Produkt der zweiten und dritten, dem doppelten Produkt der ersten und dritten sowie dem doppelten Produkt der ersten und zweiten. Schließlich werden wir die bestimmten ganzen Zahlen, mit denen diese Quadrate und doppelten Produkte multipliziert sind, in derselben Reihenfolge den \textit{ersten, zweiten, dritten, vierten, fünften, sechsten Koeffizienten} nennen. So ist
\[
axx+a'x'x'+a''x''x''+2bx'x''+2b'xx''+2b''xx'
\]
eine richtig angeordnete ternäre Form, deren erste Unbestimmte $x$, zweite $x'$, dritte $x''$, erster Koeffizient $a$ usw., vierter $b$ usw.\ ist. Es wird aber die Schreibweise sehr verkürzen, wenn wir die Unbestimmten einer ternären Form nicht immer durch besondere Buchstaben bezeichnen müssen. Deshalb werden wir dieselbe Form, wenn wir von den Unbestimmten absehen, auch so
\[
\begin{pmatrix}
a, & a', & a'' \\
b, & b', & b''
\end{pmatrix}
\]
bezeichnen.
\SourceLeaf{306}

Setzt man
\[
\begin{array}{lll}
bb-a'a''=A, & b'b'-aa''=A', & b''b''-aa'=A'' \\
ab-b'b''=B, & a'b'-bb''=B', & a''b''-bb'=B''
\end{array}
\]
so entsteht eine andere Form
\[
\begin{pmatrix}
A, & A', & A'' \\
B, & B', & B''
\end{pmatrix}\ldots F
\]
die wir die \textit{der Form}
\[
\begin{pmatrix}
a, & a', & a'' \\
b, & b', & b''
\end{pmatrix}\ldots f
\]
\textit{adjungierte} nennen werden. Bezeichnet man der Kürze halber die Zahl
\[
abb+a'b'b'+a''b''b''-aa'a''-2bb'b''\quad\text{mit}\quad D
\]
so findet man wiederum
\[
\begin{array}{lll}
BB-A'A''=aD, & B'B'-AA''=a'D, & B''B''-AA'=a''D \\
AB-B'B''=bD, & A'B'-BB''=b'D, & A''B''-BB'=b''D
\end{array}
\]
Daraus folgt, dass die der Form $F$ adjungierte Form die Form
\[
\begin{pmatrix}
aD, & a'D, & a''D \\
bD, & b'D, & b''D
\end{pmatrix}
\]
ist. Die Zahl $D$, von deren Beschaffenheit die Eigenschaften der ternären Form $f$ vor allem abhängen, werden wir die \textit{Determinante} dieser Form nennen. Auf diese Weise wird die Determinante der Form $F$ $=DD$, also gleich dem Quadrat der Determinante der Form $f$, der sie adjungiert ist.

So ist \textit{zum Beispiel} der ternären Form $\left(\begin{smallmatrix}29,\allowbreak &13,\allowbreak &9\\7,\allowbreak &-1,\allowbreak &14\end{smallmatrix}\right)$ die Form $\left(\begin{smallmatrix}-68,\allowbreak &-260,\allowbreak &-181\\217,\allowbreak &-111,\allowbreak &133\end{smallmatrix}\right)$ adjungiert; die Determinante beider $=1$.

Ternäre Formen mit der Determinante $0$ werden aus der folgenden Untersuchung vollständig ausgeschlossen. Denn wie in der bei anderer Gelegenheit ausführlicher zu behandelnden Theorie der ternären Formen gezeigt werden wird, sind sie nur \textit{dem Anschein nach} ternär und in Wirklichkeit binären Formen gleichwertig.

\Needspace{8\baselineskip}
\Article{268}
Wird eine ternäre Form $f$ mit der Determinante $D$ und den Unbestimmten $x$, $x'$, $x''$ (die erste $=x$ usw.) in eine ternäre Form $g$ mit der Determinante $E$ und den Unbestimmten $y$, $y'$, $y''$ durch eine Substitution der Art
\SourceLeaf{307}
\[
\begin{aligned}
x&=\alpha y+\beta y'+\gamma y'' \\
x'&=\alpha'y+\beta'y'+\gamma'y'' \\
x''&=\alpha''y+\beta''y'+\gamma''y''
\end{aligned}
\]
verwandelt, wobei alle neun Koeffizienten $\alpha$, $\beta$ usw.\ als ganze Zahlen vorausgesetzt werden, so werden wir der Kürze halber die Unbestimmten weglassen und einfach sagen, $f$ gehe in $g$ durch die Substitution
$(S)$
\[
\begin{array}{ccc}
\alpha, & \beta, & \gamma \\
\alpha', & \beta', & \gamma' \\
\alpha'', & \beta'', & \gamma''
\end{array}
\]
über, und $f$ \textit{enthalte} $g$, oder $g$ \textit{sei in $f$ enthalten}. Aus dieser Voraussetzung folgen unmittelbar sechs Gleichungen für die sechs Koeffizienten in $g$, die wir nicht aufschreiben müssen. Eine leichte Rechnung liefert daraus die folgenden Schlüsse:

I. Bezeichnet man der Kürze halber die Zahl
\[
\alpha\beta'\gamma''+\beta\gamma'\alpha''+\gamma\alpha'\beta''
-\gamma\beta'\alpha''-\alpha\gamma'\beta''-\beta\alpha'\gamma''\quad\text{mit }k
\]
so findet man nach den nötigen Umformungen $E=kkD$. Also teilt $D$ die Zahl $E$, und der Quotient ist ein Quadrat. Die Zahl $k$ spielt somit bei den Transformationen ternärer Formen eine ähnliche Rolle wie die Zahl $\alpha\delta-\beta\gamma$ in Artikel 157 bei den Transformationen binärer Formen: Sie ist eine Quadratwurzel aus dem Quotienten der Determinanten. Daher könnten wir vermuten, dass der Unterschied im Vorzeichen von $k$ auch hier einen wesentlichen Unterschied zwischen eigentlichen und uneigentlichen Transformationen und Enthaltungen begründet. Bei näherer Betrachtung ist aber klar. $f$ geht in $g$ auch durch die Substitution
\[
\begin{array}{ccc}
-\alpha, & -\beta, & -\gamma \\
-\alpha', & -\beta', & -\gamma' \\
-\alpha'', & -\beta'', & -\gamma''
\end{array}
\]
über; setzt man im Wert von $k$ für $\alpha$ den Wert $-\alpha$, für $\beta$ den Wert $-\beta$ usw., so erhält man $-k$. Diese Substitution wäre also der Substitution $S$ unähnlich, und jede ternäre Form, die eine andere auf die eine Art enthält, würde sie auch auf die andere Art enthalten. Eine solche \SourceLeaf{308}Unterscheidung wird deshalb hier ganz ausgeschlossen, da sie bei ternären Formen keinen Nutzen hat.

II. Bezeichnet man mit $F$, $G$ die den Formen $f$, $g$ jeweils adjungierten Formen, so werden die Koeffizienten in $F$ durch die Koeffizienten in $f$ bestimmt, und die Koeffizienten in $G$ durch die Werte der Koeffizienten der Form $g$, die aus den Gleichungen der Substitution $S$ bekannt sind. Drückt man die Koeffizienten der Form $f$ durch Buchstaben aus und vergleicht die Werte der Koeffizienten der Formen $F$, $G$, so bestätigt sich ohne Mühe, dass $F$ die Form $G$ enthält und durch die Substitution
$(S')$
\[
\begin{array}{ccc}
\beta'\gamma''-\beta''\gamma', & \gamma'\alpha''-\gamma''\alpha', & \alpha'\beta''-\alpha''\beta' \\
\beta''\gamma-\beta\gamma'', & \gamma''\alpha-\gamma\alpha'', & \alpha''\beta-\alpha\beta'' \\
\beta\gamma'-\beta'\gamma, & \gamma\alpha'-\gamma'\alpha, & \alpha\beta'-\alpha'\beta
\end{array}
\]
in sie verwandelt wird. Die Rechnung selbst, die keinerlei Schwierigkeiten bietet, lassen wir weg.

III. Die Form $g$ wird durch die Substitution $(S'')$
\[
\begin{array}{ccc}
\beta'\gamma''-\beta''\gamma', & \beta''\gamma-\beta\gamma'', & \beta\gamma'-\beta'\gamma \\
\gamma'\alpha''-\gamma''\alpha', & \gamma''\alpha-\gamma\alpha'', & \gamma\alpha'-\gamma'\alpha \\
\alpha'\beta''-\alpha''\beta', & \alpha''\beta-\alpha\beta'', & \alpha\beta'-\alpha'\beta
\end{array}
\]
offenbar in dieselbe Form verwandelt, in die $f$ durch
\[
\begin{array}{ccc}
k, & 0, & 0 \\
0, & k, & 0 \\
0, & 0, & k
\end{array}
\]
übergeht, also in die Form, die man erhält, indem man jeden Koeffizienten der Form $f$ mit $kk$ multipliziert. Diese Form werden wir mit $f'$ bezeichnen.

IV. Auf genau dieselbe Weise beweist man, dass die Form $G$ durch die Substitution $(S''')$
\[
\begin{array}{ccc}
\alpha, & \alpha', & \alpha'' \\
\beta, & \beta', & \beta'' \\
\gamma, & \gamma', & \gamma''
\end{array}
\]
\SourceLeaf{309}
in die Form übergeht, die aus $F$ entsteht, indem man jeden Koeffizienten mit $kk$ multipliziert. Diese Form werden wir mit $F'$ bezeichnen.

Wir werden sagen, die Substitution $S'''$ entstehe \textit{durch Transposition} der Substitution $S$. Dann entsteht offenbar auch $S$ wieder durch Transposition der Substitution $S'''$, und $S'$, $S''$ entstehen jeweils durch Transposition der anderen. --- Die Substitution $S'$ kann zweckmäßig als die der Substitution $S$ \textit{adjungierte} bezeichnet werden; somit ist der Substitution $S'''$ die Substitution $S''$ adjungiert.

\Needspace{8\baselineskip}
\Article{269}
Wenn nicht nur die Form $f$ die Form $g$ enthält, sondern auch diese jene, werden die Formen $f$, $g$ \textit{äquivalent} genannt. In diesem Fall teilt also nicht nur $D$ die Zahl $E$, sondern auch $E$ die Zahl $D$, woraus sich leicht ergibt, dass $D=E$ gelten muss. Enthält umgekehrt eine Form $f$ eine Form $g$ mit derselben Determinante, so sind diese beiden Formen äquivalent. Denn mit denselben Bezeichnungen wie im vorigen Artikel und unter Ausschluss des Falls $D=0$ gilt $k=\pm1$; somit ist die Form $f'$. in die $g$ durch die Substitution $S''$ übergeht, mit $f$ identisch, oder $f$ ist in $g$ enthalten. Ferner sind in diesem Fall offenbar auch die Formen $F$, $G$, die den Formen $f$, $g$ adjungiert sind, einander äquivalent, und die letztere geht durch die Substitution $S'''$ in die erstere über. Werden schließlich umgekehrt die Formen $F$, $G$ als äquivalent \textit{vorausgesetzt} und geht die erstere durch die Substitution $T$ in die letztere über, so sind auch die Formen $f$, $g$ äquivalent: $f$ geht in $g$ durch die der Substitution $T$ adjungierte Substitution über, und $g$ in $f$ durch die Substitution, die durch Transposition der Substitution $T$ entsteht. Denn durch diese beiden Substitutionen geht jeweils die der Form $F$ adjungierte Form in die der Form $G$ adjungierte Form und diese in jene über. Diese beiden Formen entstehen aber aus $f$, $g$, indem man jeden Koeffizienten mit $D$ multipliziert; daraus folgt ohne Mühe, dass durch dieselben Substitutionen jeweils $f$ in $g$. und $g$ in $f$ übergeht.

\Needspace{8\baselineskip}
\Article{270}
Wenn eine ternäre Form $f$ eine ternäre Form $f'$ enthält und diese eine Form $f''$: dann enthält auch $f$ die Form $f''$. Denn es ist sehr leicht zu sehen, dass, wenn die folgenden Übergänge stattfinden:
\[
\begin{array}{c|c}
f\text{ in }f'\text{ durch Substitution} & f'\text{ in }f''\text{ durch Substitution} \\[2pt]
\begin{array}{ccc}
\alpha, & \beta, & \gamma \\
\alpha', & \beta', & \gamma' \\
\alpha'', & \beta'', & \gamma''
\end{array}
&
\begin{array}{ccc}
\delta, & \varepsilon, & \zeta \\
\delta', & \varepsilon', & \zeta' \\
\delta'', & \varepsilon'', & \zeta''
\end{array}
\end{array}
\]
\SourceLeaf{310}
$f$ durch die Substitution
\[
\begin{array}{ccc}
\alpha\delta+\beta\delta'+\gamma\delta'', & \alpha\varepsilon+\beta\varepsilon'+\gamma\varepsilon'', & \alpha\zeta+\beta\zeta'+\gamma\zeta'' \\
\alpha'\delta+\beta'\delta'+\gamma'\delta'', & \alpha'\varepsilon+\beta'\varepsilon'+\gamma'\varepsilon'', & \alpha'\zeta+\beta'\zeta'+\gamma'\zeta'' \\
\alpha''\delta+\beta''\delta'+\gamma''\delta'', & \alpha''\varepsilon+\beta''\varepsilon'+\gamma''\varepsilon'', & \alpha''\zeta+\beta''\zeta'+\gamma''\zeta''
\end{array}
\]
verwandelt wird.

In dem Fall also, dass $f$ zu $f'$ und $f'$ zu $f''$ äquivalent ist, ist auch die Form $f$ zur Form $f''$ äquivalent. --- Wie diese Sätze auf mehrere Formen anzuwenden sind, ist im Übrigen unmittelbar klar.
